\documentclass[10pt,a4paper]{amsart}
\usepackage{amsmath,amssymb,amsthm,graphicx}
\usepackage{verbatim}
\newenvironment{reference}{\comment}{\endcomment}
%\newenvironment{reference}{}{}
\newenvironment{slogan}{\comment}{\endcomment}
\newenvironment{history}{\comment}{\endcomment}
% For commutative diagrams we use Xy-pic
\usepackage[all]{xy}

% We use 2cell for 2-commutative diagrams.
\xyoption{2cell}
% Theorem environments.
%
\theoremstyle{plain}
\newtheorem{theorem}[subsection]{Theorem}
\newtheorem{proposition}[subsection]{Proposition}
\newtheorem{lemma}[subsection]{Lemma}

\theoremstyle{definition}
\newtheorem{definition}[subsection]{Definition}
\newtheorem{example}[subsection]{Example}
\newtheorem{exercise}[subsection]{Exercise}
\newtheorem{situation}[subsection]{Situation}

\theoremstyle{remark}
\newtheorem{remark}[subsection]{Remark}
\newtheorem{remarks}[subsection]{Remarks}

\numberwithin{equation}{subsection}

% Macros
%
\def\lim{\mathop{\mathrm{lim}}\nolimits}
\def\colim{\mathop{\mathrm{colim}}\nolimits}
\def\Spec{\mathop{\mathrm{Spec}}}
\def\Hom{\mathop{\mathrm{Hom}}\nolimits}
\def\Ext{\mathop{\mathrm{Ext}}\nolimits}
\def\SheafHom{\mathop{\mathcal{H}\!\mathit{om}}\nolimits}
\def\SheafExt{\mathop{\mathcal{E}\!\mathit{xt}}\nolimits}
\def\Sch{\mathit{Sch}}
\def\Mor{\mathop{\mathrm{Mor}}\nolimits}
\def\Ob{\mathop{\mathrm{Ob}}\nolimits}
\def\Sh{\mathop{\mathit{Sh}}\nolimits}
\def\NL{\mathop{N\!L}\nolimits}
\def\CH{\mathop{\mathrm{CH}}\nolimits}
\def\proetale{{pro\text{-}\acute{e}tale}}
\def\etale{{\acute{e}tale}}
\def\QCoh{\mathit{QCoh}}
\def\Ker{\mathop{\mathrm{Ker}}}
\def\Im{\mathop{\mathrm{Im}}}
\def\Coker{\mathop{\mathrm{Coker}}}
\def\Coim{\mathop{\mathrm{Coim}}}

% Boxtimes
%
\DeclareMathSymbol{\boxtimes}{\mathbin}{AMSa}{"02}

%
% Macros for moduli stacks/spaces
%
\def\QCohstack{\mathcal{QC}\!\mathit{oh}}
\def\Cohstack{\mathcal{C}\!\mathit{oh}}
\def\Spacesstack{\mathcal{S}\!\mathit{paces}}
\def\Quotfunctor{\mathrm{Quot}}
\def\Hilbfunctor{\mathrm{Hilb}}
\def\Curvesstack{\mathcal{C}\!\mathit{urves}}
\def\Polarizedstack{\mathcal{P}\!\mathit{olarized}}
\def\Complexesstack{\mathcal{C}\!\mathit{omplexes}}
% \Pic is the operator that assigns to X its picard group, usage \Pic(X)
% \Picardstack_{X/B} denotes the Picard stack of X over B
% \Picardfunctor_{X/B} denotes the Picard functor of X over B
\def\Pic{\mathop{\mathrm{Pic}}\nolimits}
\def\Picardstack{\mathcal{P}\!\mathit{ic}}
\def\Picardfunctor{\mathrm{Pic}}
\def\Deformationcategory{\mathcal{D}\!\mathit{ef}}
\usepackage{cleveref}
\crefname{originalreference}{original reference}{original references}
\Crefname{originalreference}{Original reference}{Original references}
\makeatletter
\newlabel{spaces-descent-lemma-descending-property-isomorphism}{{spaces-descent-lemma-descending-property-isomorphism}{source}{}{sourcepointer-1}{}}
\newlabel{spaces-descent-lemma-descending-property-isomorphism@cref}{{[originalreference][0][]spaces-descent-lemma-descending-property-isomorphism}{[1][0][]source}{}{}{}}
\makeatother
\title{The property that a morphism induces an isomorphism of automorphism groups is independent of the chosen point in its class}
\author{The Stacks Project Authors}
\date{stacks-project:stacks-morphisms.tex:local-snapshot-20261012; GNU FDL 1.2 or later}
\begin{document}
\maketitle
\noindent\textit{Editorial note.} All delivered mathematical text is the author's own, copied from the pinned Stacks Project chapter source sources/stacks-algebraic-stacks-morphisms/pinned-original-source.tex; no statement or proof step was rewritten, reordered or removed.\par\medskip
\noindent\textit{Editorial note.} Preamble provenance: the Stacks Project keeps one single shared preamble for every chapter; this card takes it directly from the staged sources/stacks-algebraic-stacks-morphisms/preamble.tex as three verbatim spans (lines 8-12, 14-18 and 154-217), so no derived file was written and nothing was added or redefined. The wrapper class and the base packages are supplied by the assembler.\par\medskip
\noindent\textit{Editorial note.} Bibliography provenance: neither the counted statement nor its proof contains a citation (this lane rescans the whole span before writing the card), so the card carries no bibliography and no entry was reconstructed.\par\medskip
\noindent\textit{Editorial note.} Selection provenance: this card was opened under the file-based control contract (CONTROL.md), with the mechanical pre-card gate run on its statement and proof interval before the card was written.\par\medskip
\noindent\textit{Allowed standard prerequisites (editorial).}\begin{itemize}
\item The automorphism group algebraic spaces of a point of an algebraic stack, and the equivalence relation on such points given by field extensions. The author's section on automorphism groups, in which the counted statement occurs. Fixes the two maps whose being isomorphisms is asserted to be equivalent.
\item The descent of the property of being an isomorphism along a faithfully flat and quasi-compact morphism. Invoked at the start of the counted proof (source lines 4030-4033) as a cross-chapter pointer. Shows that the property propagates along a field extension, whence along the equivalence class of the point.
\end{itemize}\medskip
\ifdefined\hypertarget\hypertarget{sourcepointer-1}{}\fi
\noindent\textit{Source reference.} Original reference spaces-descent-lemma-descending-property-isomorphism: sources/stacks-algebraic-stacks-morphisms/pinned-original-source.tex, source lines 4033--4033. Navigation pointer only: the counted text refers to the Stacks item spaces-descent-lemma-descending-property-isomorphism, which belongs to another chapter of the project and is not part of this card. The pointer carries the target's own label as its number because this chapter's staged snapshot does not contain the other chapter's tag table; its content is not reproduced here.\par\medskip
\begin{lemma}
\label{lemma-iso-automorphism-groups}
Let $f : \mathcal{X} \to \mathcal{Y}$ be a morphism of algebraic stacks.
Let $x \in |\mathcal{X}|$ be a point. The following are equivalent
\begin{enumerate}
\item for some morphism $x : \Spec(k) \to \mathcal{X}$ in the
class of $x$ setting $y = f \circ x$ the map
$G_x \to G_y$ of automorphism group algebraic spaces
is an isomorphism, and
\item for any morphism $x : \Spec(k) \to \mathcal{X}$ in the
class of $x$ setting $y = f \circ x$ the map
$G_x \to G_y$ of automorphism group algebraic spaces
is an isomorphism.
\end{enumerate}
\end{lemma}

\begin{proof}
This comes down to the fact that being an isomorphism
is fpqc local on the target, see
Descent on Spaces, Lemma
\ref{spaces-descent-lemma-descending-property-isomorphism}.
Namely, suppose that $k'/k$ is an extension of fields and
denote $x' : \Spec(k') \to \mathcal{X}$ the composition
and set $y' = f \circ x'$.
Then the morphism $G_{x'} \to G_{y'}$ is the base change
of $G_x \to G_y$ by $\Spec(k') \to \Spec(k)$.
Hence $G_x \to G_y$ is an isomorphism
if and only if $G_{x'} \to G_{y'}$ is an isomorphism.
Thus we see that the property propagates through the
equivalence class if it holds for one.
\end{proof}

\end{document}
